%% ma: L10-B14-0019
%% lop: 10
%% chuong: 5
%% bai: 14
%% dang: 10.14.4
%% loai: TN4
%% muc: VD
%% dap_an: "C"
%% nguon: "chuyen-de-cac-so-dac-trung-cua-mau-so-lieu-khong-ghep-nhom-toan-10-chuong-trinh-moi.pdf (GV. Phan Nhật Linh) (P04), Bài 03 Dạng 1 Phần I Câu 9"
%% kien_thuc: "Lớp 10 Bài 13 (số trung bình, tứ phân vị), Bài 14 (khoảng biến thiên, khoảng tứ phân vị, giá trị bất thường)"
%% trang_thai: chinh-thuc
%% ghi_chu: "Lời giải gốc ghi nhầm trung bình 355/7 (đúng 335/7). Tự duyệt theo yêu cầu giáo viên 22/10/2026."
\begin{ex}
Điều tra năng suất lúa (tạ/ha) của $7$ héc-ta trồng lúa ở hai vùng A và B, ta thu được mẫu số liệu sau:
\begin{center}
\begin{tabular}{|c|c|c|c|c|c|c|c|}\hline
Vùng A & $41$ & $44$ & $45$ & $47$ & $51$ & $53$ & $54$\\\hline
Vùng B & $43$ & $44$ & $47$ & $48$ & $50$ & $51$ & $52$\\\hline
\end{tabular}
\end{center}
Khẳng định nào dưới đây là \textbf{sai}?
\choice{Năng suất trung bình của hai vùng A và B như nhau}{Khoảng tứ phân vị của vùng A lớn hơn khoảng tứ phân vị của vùng B}{\True Vùng A trồng lúa ổn định hơn vùng B}{Khoảng biến thiên của vùng A lớn hơn khoảng biến thiên của vùng B}
\loigiai{Năng suất trung bình của cả hai vùng đều bằng $\dfrac{335}7$: A đúng.

Vùng A: $Q_1=44$, $Q_3=53$, $\Delta_Q=9$; vùng B: $Q_1=44$, $Q_3=51$, $\Delta_Q=7$: B đúng.

$R_A=54-41=13>R_B=52-43=9$: D đúng; vùng B ổn định hơn nên C sai. Chọn C.}
\end{ex}
